% ═══════════════════════════════════════════════════════════════════════
%  Stratégie 2 --- XAUUSD niveaux x10. Chiffres : macros et tables générées.
% ═══════════════════════════════════════════════════════════════════════
\section{Robustesse avancée et table de verdict}
\label{sec:x10_robustesse}

\lettrine[lines=2, findent=2pt, nindent=0pt]{L}{es} tests de la section
précédente répondaient à « le résultat est-il stable~? ». Ceux-ci répondent à
une question différente et plus sévère~: \emph{ce résultat aurait-il pu être
produit par le hasard, compte tenu du nombre de fois où l'on a cherché~?}
L'annexe~\ref{app:x10_methodologie_robustesse} explique chaque test en
détail~; cette section publie les mesures.

% GÉNÉRÉ par scripts/build_x10_report_assets.py — ne pas éditer à la main.
% Toute valeur vient de results/xau_x10/*.json.
\begin{table}[htbp]
\centering
\footnotesize
\caption{Batterie de robustesse de la configuration analysée. Les intervalles sont des bootstraps stationnaires par blocs~; le DSR et le haircut sont déflatés par le registre d'essais.}
\label{tab:x10_robustness}
\begin{tabular}{l r l r}
\toprule
\textbf{Mesure} & \textbf{Valeur} & \textbf{IC 95\,\% / repère} & \textbf{Seuil} \\
\midrule
Sharpe net (séance, 252) & $-1{,}942$ & [$-2{,}686$ ; $-1{,}285$] & --- \\
PSR (contre $0$) & $0{,}000$ & --- & --- \\
DSR (déflaté, $N = 34$) & $0{,}000$ & --- & $\geq 0{,}95$ \\
Sharpe attendu du maximum de $N$ essais & $0{,}636$ & --- & --- \\
Haircut Sharpe (BHY) & $-0{,}194$ & --- & $p$ ajustée $1{,}000$ \\
MinBTL & \emph{non calculable} & --- & --- \\
PBO (CSCV, $16$ bins) & $0{,}2460$ & --- & $\leq 0{,}35$ \\
\midrule
Rendement total & $-67{,}2$\,\% & [$-77{,}8$\,\% ; $-53{,}0$\,\%] & --- \\
Repli maximal & $68{,}0$\,\% & [$54{,}9$\,\% ; $78{,}2$\,\%] & --- \\
Profit factor & $0{,}673$ & [$0{,}581$ ; $0{,}766$] & --- \\
Ratio de Sortino & $-2{,}843$ & [$-3{,}735$ ; $-1{,}968$] & --- \\
\midrule
Repli maximal observé (Monte-Carlo) & $34{,}7$\,\% & médiane $34{,}6$\,\% ; p95 $36{,}4$\,\% & --- \\
$z$ de chance de séquence & $0{,}14$ & --- & --- \\
P(ruine) & $0{,}0$\,\% & --- & --- \\
P(perte $> 50$\,\%) & $99{,}3$\,\% & --- & --- \\
Capital terminal médian (bootstrap) & $0{,}326$ & --- & --- \\
\bottomrule
\end{tabular}
\end{table}


\subsection{Deflated Sharpe Ratio}

Chercher un bon réglage, c'est tirer plusieurs fois. Plus on tire, plus le
meilleur tirage paraît bon --- même si tous les tirages sont du bruit. Le
\emph{Deflated Sharpe Ratio} corrige cet effet~: il compare le Sharpe observé
non pas à zéro, mais au Sharpe qu'on attendrait du \emph{meilleur} de $N$
essais tirés au hasard.

\begin{insightbox}
\textbf{Avec \XTenTrialsCount{} essais, le meilleur d'entre eux aurait un
Sharpe attendu de \metric{\XTenExpectedMaxSharpe} par pur hasard.} C'est le
seuil que la stratégie devait franchir, et il est positif. Le Sharpe observé
vaut \XTenSharpeIS~: le \code{DSR} vaut donc \metric{\XTenDSR}, contre le
seuil gelé de $0{,}95$. Le \emph{Probabilistic Sharpe Ratio} --- la même
mesure sans déflation, contre zéro --- vaut \XTenPSR. Les deux sont nuls à la
précision publiée.
\end{insightbox}

Deux mesures voisines complètent le tableau. Le \emph{haircut Sharpe} avec
correction pour tests multiples ramène le Sharpe à \XTenHaircutSharpe{} pour
une $p$-valeur ajustée de \XTenHaircutPValue~: autrement dit, le résultat est
strictement indistinguable de l'absence de signal. Et la \emph{longueur
minimale de backtest} n'est pas calculable~: elle est définie pour un Sharpe
observé positif, ce qui n'est pas le cas ici.

\subsection{Probability of Backtest Overfitting}

Le \code{PBO} pose une question empirique et intuitive~: si l'on choisit la
meilleure configuration sur une moitié des données, quelle est la probabilité
qu'elle se classe dans la moitié \emph{inférieure} sur l'autre~? Il se mesure
en découpant l'échantillon en blocs et en examinant toutes les manières de
les répartir --- ici \XTenPBOSplits{} combinaisons sur les vingt-sept
configurations.

Le \code{PBO} mesuré vaut \metric{\XTenPBO}, contre un seuil gelé de
$0{,}35$. Le critère échoue.

\begin{warningbox}
\textbf{Attention à ce que ce chiffre dit, et à ce qu'il ne dit pas.} Un
\code{PBO} élevé sur une grille qui \emph{perd partout} ne signifie pas « le
surajustement explique la perte ». Il signifie que le classement des
vingt-sept configurations entre les deux moitiés de l'échantillon n'a presque
aucune persistance --- ce qui est exactement ce à quoi on s'attend quand il
n'y a rien à classer. \textbf{Il n'y a pas de surajustement parce qu'il n'y a
pas de signal à surajuster.} Le critère échoue, et la raison de son échec est
cohérente avec tout le reste du dossier.
\end{warningbox}

\subsection{Validation croisée combinatoire purgée}

La validation croisée combinatoire purgée découpe l'échantillon en six
groupes, en teste deux à la fois, et purge les observations à cheval sur la
frontière pour éviter qu'une information ne fuite d'un côté à l'autre. Elle
produit \XTenCPCVSplits{} découpages et plusieurs chemins reconstruits, sur
les vingt-sept configurations.

% GÉNÉRÉ par scripts/build_x10_report_assets.py — ne pas éditer à la main.
% Toute valeur vient de results/xau_x10/*.json.
\begin{table}[htbp]
\centering
\small
\caption{Validation croisée combinatoire purgée~: les cinq meilleures configurations sur la médiane hors échantillon. Aucune des vingt-sept n'est positive sur un seul des splits. La métrique est le Sharpe des rendements minute du portefeuille et ne se compare au reste du rapport qu'en signe.}
\label{tab:x10_cpcv}
\begin{tabular}{l r r r r}
\toprule
\textbf{Configuration} & \textbf{Médiane OOS} & \textbf{q05} & \textbf{q95} & \textbf{Splits positifs} \\
\midrule
$z=0{,}5$ ; $a_{\min}=0{,}3$ ; $k_s=1{,}50$ & $-0{,}959$ & $-1{,}515$ & $-0{,}447$ & $0$\,\% \\
$z=0{,}5$ ; $a_{\min}=0{,}3$ ; $k_s=0{,}75$ & $-0{,}975$ & $-1{,}662$ & $-0{,}274$ & $0$\,\% \\
$z=1{,}5$ ; $a_{\min}=0{,}3$ ; $k_s=0{,}75$ & $-0{,}996$ & $-1{,}882$ & $-0{,}533$ & $0$\,\% \\
$z=0{,}5$ ; $a_{\min}=0{,}3$ ; $k_s=1{,}00$ & $-1{,}007$ & $-1{,}563$ & $-0{,}390$ & $0$\,\% \\
$z=1{,}5$ ; $a_{\min}=0{,}3$ ; $k_s=1{,}50$ & $-1{,}015$ & $-1{,}609$ & $-0{,}776$ & $0$\,\% \\
\bottomrule
\end{tabular}
\end{table}


\textbf{Aucune des vingt-sept configurations n'a une médiane hors échantillon
positive}, et aucune n'est positive sur un seul des découpages. La meilleure
médiane vaut \XTenCPCVBestMedian, la moyenne sur toute la grille
\XTenCPCVMean, et la part de configurations à médiane positive vaut
\XTenCPCVPositiveShare.

La métrique utilisée ici n'est pas celle de la grille de la
section~\ref{sec:x10_in_sample}~: c'est un Sharpe calculé sur les rendements
minute du portefeuille, convention du module de validation croisée. Les deux
ne se comparent qu'en \emph{signe}, jamais en niveau, et ce rapport ne les met
pas en regard.

\subsection{Deux tests dont les résultats ne sont pas publiables}

L'honnêteté impose de signaler ce qui n'a pas fonctionné, plutôt que de
l'omettre.

\begin{description}
  \item[Le test de Romano--Wolf a échoué.] Sa procédure descendante se
  retrouve avec un tableau vide lorsque aucune configuration de la grille ne
  bat le portefeuille plat. C'est une limite préexistante du module de tests
  statistiques du dépôt, constatée et \emph{non corrigée} par cette campagne.
  \item[Le test de supériorité prédictive n'est pas interprétable en
  l'état.] Le module du dépôt transmet des \emph{rendements} à une
  implémentation qui documente attendre des \emph{pertes}. L'orientation est
  donc inversée, et la $p$-valeur rendue signifie l'inverse de ce qu'elle
  paraît dire sur une grille intégralement perdante. Le chiffre existe dans
  les fichiers de résultats, assorti de son avertissement~; \textbf{il n'est
  pas publié ici}, et aucun critère de la table de décision n'en dépend.
\end{description}

\begin{insightbox}
\textbf{Un test cassé qui semble donner raison est plus dangereux qu'un test
absent.} Publié tel quel, celui-ci aurait fait passer une grille entièrement
perdante pour une grille statistiquement significative. Il a été identifié,
documenté dans les résultats, et écarté du rapport. Sa correction relève du
module de tests, pas de cette campagne.
\end{insightbox}

\subsection{Table de verdict}

Les neuf critères gelés avant la mesure, confrontés aux résultats, donnent la
table reproduite au résumé (section~\ref{sec:x10_resume_verdict}) et en
annexe~\ref{app:x10_table_decision}. Son résultat tient en trois lignes~:

\begin{itemize}
  \item \XTenCriteriaPassing{} critère sur \XTenCriteriaCount{} est franchi~:
  celui du nombre de transactions ;
  \item \XTenCriteriaFailing{} sont en échec, dont six avec une mesure
  interprétable ;
  \item deux le sont par \emph{non-mesure} --- le portage MetaTrader~5 et la
  lecture hors échantillon --- et la table de décision impose de les compter
  comme des échecs, jamais comme « non applicables ».
\end{itemize}

\ifdefstring{\XTenVerdict}{NOGO}{%
\begin{warningbox}[title=\textbf{\color{white}\textsf{Conséquence mécanique}}]
La table de décision énonce~: « au moins un critère en échec avec une mesure
interprétable » $\Rightarrow$ \textbf{ne pas déployer}. Il y en a six. La
section~\ref{sec:x10_non_deploiement} en tire les conséquences~; la section
« conditions de déploiement », prévue sous verdict favorable, n'est pas
livrée, parce qu'elle n'a pas lieu d'être.
\end{warningbox}
}{}
